% rollout_feedback.tex
% A compact, paper-friendly LaTeX rendering of one self-learning feedback report.
% You can either compile this file standalone or copy the body into your appendix.



\hypersetup{
    colorlinks=true,
    linkcolor=blue!55!black,
    urlcolor=blue!55!black,
    citecolor=blue!55!black
}

% ---------- Styles ----------
\tcbset{
  rolloutsummary/.style={
    enhanced,
    breakable,
    colback=gray!2,
    colframe=black!25,
    boxrule=0.35pt,
    arc=1.5pt,
    left=6pt,
    right=6pt,
    top=5pt,
    bottom=5pt,
    before skip=6pt,
    after skip=6pt,
  },
  rolloutfailure/.style={
    enhanced,
    breakable,
    colback=gray!1,
    colframe=black!22,
    boxrule=0.35pt,
    arc=1.5pt,
    left=6pt,
    right=6pt,
    top=5pt,
    bottom=5pt,
    before skip=6pt,
    after skip=6pt,
  },
  rolloutnote/.style={
    enhanced,
    breakable,
    colback=blue!2,
    colframe=blue!18,
    boxrule=0.35pt,
    arc=1.5pt,
    left=6pt,
    right=6pt,
    top=5pt,
    bottom=5pt,
    before skip=6pt,
    after skip=6pt,
  }
}

\newcommand{\nodeid}[1]{\texttt{#1}}
\newcommand{\metric}[1]{\textbf{#1}}
\newcommand{\failtag}[1]{\texttt{\textbf{#1}}}
\newcommand{\code}[1]{\texttt{#1}}

\begin{tcolorbox}[rolloutnote]
\small
\textbf{Summary.}
In the first rehearsal iteration, perception succeeds across all sampled environments, while downstream failures arise from grasp instability and incorrect pan placement. The placement node fails because the pan is released away from the burner footprint, yielding zero burner coverage.
\end{tcolorbox}

\begin{tcolorbox}[rolloutsummary]
\small
\textbf{Iteration 1 Feedback} \hfill \textbf{Terminal coverage: 0.00}

\vspace{4pt}
\renewcommand{\arraystretch}{1.18}
\begin{tabularx}{\linewidth}{@{}p{0.23\linewidth}p{0.18\linewidth}X@{}}
\toprule
\textbf{Node} & \textbf{Validation} & \textbf{Feedback} \\
\midrule
\nodeid{perceive\_pan}
& \metric{1.00 (30/30)}
& Perception publishes the target OBB for the pan handle. \\

\nodeid{perceive\_burner}
& \metric{1.00 (30/30)}
& Perception publishes the container OBB for the burner. \\

\nodeid{grasp\_pan}
& \metric{0.83 (25/30)}
& The pan is considered grasped when it remains in contact with a robot link at the exit of the grasp subgraph. \\

\nodeid{place\_pan}
& \metric{0.57 (17/30)}
& The pan should be released by the gripper and cover at least 70\% of the burner XY footprint. All evaluated placement attempts fail because the pan does not cover the burner footprint. \\
\bottomrule
\end{tabularx}
\end{tcolorbox}

\begin{tcolorbox}[rolloutfailure]
\small
\textbf{Representative Failure Cases}

\vspace{2pt}
\begin{itemize}[leftmargin=1.2em, itemsep=3pt, topsep=2pt]
    \item \failtag{env 3, grasp\_pan}: grasp failed with the gripper still open.
    \[
    \code{gripper\_open\_fraction}=1.0,\quad
    \code{pan\_z}=0.024,\quad
    \code{contact}=[\code{table\_top}]
    \]
    The subgraph returned \code{failed}; only \code{ee\_pose\_at\_grasp} was bound. 

    \item \failtag{env 9, grasp\_pan}: grasp failed after gripper closure.
    \[
    \code{gripper\_open\_fraction}: 1.0 \rightarrow 0.0,\quad
    \code{pan\_z}=0.024,\quad
    \code{contact}=[\code{table\_top}]
    \]
    The robot closed the gripper, but the pan remained on the table and the node returned \code{failed}. The elapsed time was \code{4.377 s}.

    \item \failtag{place\_pan}: the evaluated placement failure trials has:
    \[
\code{pan\_coverage\_of\_burner}\leq 0.7
    \]
    Although the pan was released from the gripper, its projected bottom footprint did not sufficiently overlap with the burner footprint.
\end{itemize}
\end{tcolorbox}

\begin{tcolorbox}[rolloutnote]
\small
\textbf{Optimization signal.}
The feedback suggests that the next graph update should revise the placement target and release pose for \nodeid{place\_pan}, while also improving grasp robustness for the failed \nodeid{grasp\_pan} environments.
\end{tcolorbox}


